\documentclass[11pt,a4paper]{article}
\usepackage[margin=26mm]{geometry}
\usepackage{amsmath,amssymb,amsthm,hyperref}
\newtheorem{proposition}{Proposition}
\title{A hypermap whose distinguished involutions fail the intersection property}
\author{Human Compression project}\date{}
\begin{document}\maketitle
The unnamed paragraph opening Section~5.4 of Leemans's fixed survey
edition~\cite{survey} states that the hypermaps of Conder--Poto\v cnik--
\v Sir\'a\v n~\cite{CPS} satisfy the intersection property, citing the
absence of nontrivial cyclic normal subgroups and Conder--Oliveros,
Corollary~4.2. The latter is the string result reproduced as survey
Theorem~3.2: a pair of the three involutions must commute. That condition
cannot be discarded. We give an explicit counterexample to the unnamed
inference, without questioning the hypermap classification in~\cite{CPS}.

\begin{proposition}
There is a regular hypermap with full group $\operatorname{PSL}_2(11)$
and type $(6,6,5)$ whose three distinguished involutions do not form a
C-group marking.
\end{proposition}
\begin{proof}
All matrices below have entries in $\mathbb F_{11}$. Put
\[
 A=\begin{pmatrix}5&1\\7&6\end{pmatrix},\qquad
 B=\begin{pmatrix}0&1\\10&0\end{pmatrix},\qquad
 C=\begin{pmatrix}4&3\\9&7\end{pmatrix}.
\]
They have determinant one and trace zero, so their images $a,b,c$ in
$\operatorname{PSL}_2(11)$ are involutions. Direct multiplication gives
\[
 ABACBABCAC=-U,\qquad U=\begin{pmatrix}1&1\\0&1\end{pmatrix}.
\]
The matrices $U$ and $B$ generate $\operatorname{SL}_2(11)$. Indeed the
powers of $U$ give every elementary upper unipotent matrix, and conjugating
by $B$ gives every elementary lower unipotent matrix. These perform row
elimination. To include the resulting diagonal matrices explicitly, for
$s\ne0$ one has
\[
 \begin{pmatrix}1&s\\0&1\end{pmatrix}
 \begin{pmatrix}1&0\\-s^{-1}&1\end{pmatrix}
 \begin{pmatrix}1&s\\0&1\end{pmatrix}
 =\begin{pmatrix}s&0\\0&s^{-1}\end{pmatrix}B.
\]
Thus $\langle a,b,c\rangle=\operatorname{PSL}_2(11)$.

Two more direct products are
\[
 X=(AB)^3=\begin{pmatrix}3&10\\10&8\end{pmatrix},\qquad
 Z=(BC)^3=\begin{pmatrix}1&3\\3&10\end{pmatrix},\qquad XZ=-B.
\]
Write $x,z$ for their projective images. Since $x=(ab)^3$ and
$x=bz^{-1}$, the element $x$ belongs to
$\langle a,b\rangle\cap\langle b,c\rangle$. The displayed matrix $X$
is neither $\pm I$ nor $\pm B$. Therefore
\[
 \langle a,b\rangle\cap\langle b,c\rangle
                  \supsetneq\langle b\rangle,
\]
which is the required failure of the intersection property. The group
$\operatorname{PSL}_2(11)$ is nonabelian simple, so it has no nontrivial
cyclic normal subgroup. Here we retain the standard simplicity theorem
for $\operatorname{PSL}_2(q)$ with $q\ge4$; no classification search is
needed to infer simplicity from a group order.

Finally put $r=ab$ and $s=bc$. Their orders are six and six, while
$rs=ac$ has order five; these orders follow by multiplying the displayed
matrices and checking powers modulo $\{I,-I\}$. The equation $xz=b$
gives $b=r^3s^3$, so $r,s$ generate the same group as $a,b,c$.
Conjugation by the involution $b$ inverts both $r$ and $s$.
The type is hyperbolic, since $1/6+1/6+1/5<1$.
The introduction of~\cite{CPS} identifies exactly this situation as an
inner-reflexible generating triple: it gives a regular hypermap on a
nonorientable surface with full automorphism group $\langle r,s\rangle$.
This establishes the hypermap interpretation as well as the group calculation.
\end{proof}

The valid string theorem remains unchanged. Hypermap enumeration can
still provide candidates for C-group markings, but the appropriate
intersection tests must be imposed. The blanket transfer made in the
survey is false. The additional literal assertion about all values of
$q$ also needs its field range: $\operatorname{PSL}_2(2)=S_3$ has a normal
cyclic subgroup of order three. The main hyperbolic scope of~\cite{CPS}
already excludes that elliptic example. Our counterexample has $q=11$
and lies within that hyperbolic scope, so it is independent of this
small-field qualification.

The displayed certificate was found while auditing this particular
survey clause. This note does not claim to settle a named open conjecture,
to replace the full hypermap enumeration, or to complete the survey's
other numerical obligations. The hypermap correspondence and original
enumeration remain attributed to~\cite{CPS}.
\begin{thebibliography}{9}
\bibitem{survey} D.~Leemans, \emph{String C-group representations of
almost simple groups: a survey}, arXiv:1910.08843v2, 11 May 2020,
unnamed paragraph at the start of Section~5.4.
\bibitem{CPS} M.~Conder, P.~Poto\v cnik and J.~\v Sir\'a\v n,
\emph{Regular hypermaps over projective linear groups},
J.~Aust.~Math.~Soc.~85 (2008),155--175, especially pp.~156--157
and Section~6; doi:10.1017/S1446788708000827.
\end{thebibliography}
\end{document}
